\documentclass[10pt]{article}
\usepackage[margin=20mm]{geometry}
\usepackage{amsmath,amssymb,fontspec,unicode-math,xurl,hyperref}
\IfFontExistsTF{DejaVu Serif}{\setmainfont{DejaVu Serif}}{\setmainfont{Latin Modern Roman}}
\setmathfont{latinmodern-math.otf}
\setlength{\parindent}{0pt}\setlength{\parskip}{6pt}\sloppy
\begin{document}
Abstract: The reported simulated magnetization improvements are not refuted by this review. Exact Bloch and noise calculations distinguish those proxies from equilibrium recovery and concentration precision. The reduction of total measurement time at matched abundance error remains unverified from the supplied evidence.\par
Author: AI-assisted technical review\par
Date:\par
Title: Technical verification of policy-based NMR sequencing\par
\# Identity, scope and outcome\par
Rohan Shenoy, Evan A. Coleman, Hans Gaensbauer and Elsa A. Olivetti, Counting atoms faster: policy-based nuclear magnetic resonance pulse sequencing for atomic abundance measurement, ICML 2025, PMLR 267:54810–54824, OpenReview . This review covers the frozen 15-page publisher PDF, main Sections 1–4, Appendix A’s Bloch equations and Appendix B’s policy and PPO descriptions. All four result figures were actually viewed; page numbers below refer to PDF pages. An associated author repository was not identified in the frozen bounded source search. Author code was not executed.\par
The assessment is Needs Evidence. No specific numerical experiment is shown false. The results support improved simulated magnetization proxies, but do not establish a concentration-error-versus-total-time comparison at matched precision. Independent analytic checks explain why transverse suppression, peak signal and faster cycles are insufficient by themselves for that stronger inference. The paper’s own partial-observation and non-equilibrium qualifications are preserved.\par
\# What the paper demonstrates\par
Section 2.2 fits effective spin frequencies and weights to laboratory caffeine spectra and samples longitudinal relaxation parameters. Section 3 uses PPO to maximize the running peak of \(M_x^2+M_y^2\), minimize transverse magnetization, and toggle the two policies. Figure 1 reports peak magnetization against concentration with errors across five training seeds and five environments. The learned curves increase more strongly than the baseline in the plotted concentration range. Figure 2 shows a lower transverse-magnetization trace for a spoiling policy in one environment. Figures 3–4 show state-action structure and repeated magnetization cycles.\par
Section 3.1 reports approximately \(16\%\) of baseline time to attain a peak. At face value this is a time-to-peak statement, not a matched-error metrological speedup. The plotted curves and text do not specify an inverse concentration estimator, held-out concentration prediction errors, detection limits, or a total-time curve at a fixed concentration confidence interval. Five seeds and multiple simulated scans are useful evidence; they are not absent. However, their plotted uncertainty does not directly provide the required calibration error after inversion or dependence across repeated toggled scans.\par
\# Analytic checks that bound the interpretation\par
\#\# Transverse spoiling versus equilibrium recovery\par
With the paper’s Bloch equations, \(B=(b,0,B_0)\) and fixed parameters, the states \[M^+=(0,0,M_0),\qquad M^-=(0,0,-M_0)\] have identical observed transverse magnetization and identical spoiling reward zero. Nevertheless, their distances from equilibrium are \(0\) and \(2|M_0|\), and \[\dot M_y^+=\gamma bM_0,\qquad
 \dot M_y^-=-\gamma bM_0.\] For the dimensionless diagnostic \(M_0=\gamma=T_1=T_2=B_0=1\), \(b=0.2\), the derivatives are \((0,0.2,0)\) and \((0,-0.2,2)\). Thus suppressing the observed transverse component alone does not determine longitudinal recovery or the next response. This is a check of the printed equations, not a claim that the learned policy reaches \(M^-\). It identifies the unmeasured variable needed to justify accelerated equilibrium recovery.\par
Crucially, p. 6 explicitly states that the shortened spoiling interval does not reset the system to equilibrium and considers partial resets; Appendix B, p. 13, explicitly says \(M_z,M_0,T_1,T_2,\gamma\) are masked and a non-recurrent policy need not reconstruct the dynamics. Consequently, partial observability is an acknowledged scope condition, not a newly discovered fatal flaw. A stable non-equilibrium measurement cycle may be useful if calibrated. The unresolved question is its accuracy and repeatability, not whether partial-reset control is categorically invalid.\par
\#\# A peak reward is not a noise-free abundance statistic\par
In the stated Gaussian readout model, consider a simple special case with zero true transverse signal and independent noise components \(\epsilon_x,\epsilon_y\sim N(0,\sigma^2)\) at each readout. Then \(Z=\epsilon_x^2+\epsilon_y^2\) is exponential with mean \(2\sigma^2\). For \(N\) independent readouts, \[\mathbb E\max_{1\leq j\leq N} Z_j=2\sigma^2 H_N,
 \qquad H_N=\sum_{j=1}^N j^{-1}.\] This follows by integrating \(1-(1-e^{-z/(2\sigma^2)})^N\), or by exponential order-statistic spacings. The independent calculation gives \(10.3748\sigma^2\) for \(N=100\) and \(14.9709\sigma^2\) for \(N=1000\), although abundance is unchanged. The paper’s nonzero, temporally dependent, changing-temperature trajectories need not obey this special-case formula. It is not an estimate of bias in Figure 1 and does not establish that PPO exploits noise. It shows why the actual peak acquisition window and a calibrated error estimate must accompany a peak-based precision claim. If all methods use fixed matched windows and calibration absorbs the noise, this concern may be resolved.\par
\#\# Repeated scans and the square-root rule\par
Section 3.2 motivates repetitions through \(1/\sqrt n\) noise reduction. For equal-variance readouts \(X_i\), \[\operatorname{Var}(\bar X)=\frac{\sigma^2}{n}
 +\frac{2}{n^2}\sum_{i<j}\operatorname{Cov}(X_i,X_j).\] Independent measurements recover the stated rule. Equicorrelation \(\rho\) instead gives \(\sigma^2[\rho+(1-\rho)/n]\). For \(n=100\), \(\rho=0.1\) this is \(0.109\sigma^2\) rather than \(0.01\sigma^2\). This illustrative valid covariance model does not assert that the experimental covariance is positive or that electronic noise must become dependent after partial reset. Persistent physical state and independently sampled electronic noise are distinct. Actual scan covariance and estimator bias are needed to determine whether the accelerated cycle maintains the same precision per unit time.\par
\# Algorithm and source limits\par
The printed PPO clipped surrogate, GAE recursion and actor/critic roles are consistent at their stated abstract level. The simulator observes fewer variables than its full Markov state; it is best interpreted as partially observed control. At \(0.2\) ms per control step, the 200-step toggle restriction corresponds to \(40\) ms. The source does not expose the exact ODE integrator, internal tolerances/substeps, fitted spin sets, pulse implementations or historical library commit. A claim of numerical instability from the control step alone would be unsupported, since an integrator may use smaller internal steps. Hardware cost and thermal feasibility, quantum footnote details, independent laboratory validation and external baseline provenance were not certified.\par
\# Conclusion and evidence needed\par
The simulation results justify further study of the three-policy design. A stronger abundance-measurement claim needs an explicit concentration estimator evaluated on held-out concentrations, errors or confidence coverage against total acquisition time, and repeated-cycle calibration/covariance. Full-state recovery measurements are needed only for the stronger equilibrium-recovery interpretation; a calibrated partial-reset method need not restore equilibrium. No correction or failed experiment is asserted, and no later author revision was inspected.\par
Item | Assessment | Verification | Adoption\par
Bloch proxy distinctions | Scope clarification | Exact derivations | Not applicable\par
Equal-precision time reduction | Needs Evidence | Not established | Unknown\par
Historical simulation fidelity | Needs Evidence | Code unavailable | Unknown\par
1 R. Shenoy, E. A. Coleman, H. Gaensbauer and E. A. Olivetti. Counting atoms faster: policy-based nuclear magnetic resonance pulse sequencing for atomic abundance measurement. PMLR 267 (2025), 54810–54824. https://proceedings.mlr.press/v267/shenoy25a.html. Reviewed publisher version.\par
\# Reproducibility and roles\par
The matching directory contains the independent check script and its actual output, source and code hashes, coverage, findings, compiler logs and page-level visual records. Author programs were not executed. The source PDF SHA-256 is . Separate assigned GPT-6 Astra / Max units performed the initial review and the subsequent independent source-based check recorded in the verification files. Scientific adjudication remains pending; human reviews are empty.\par
\end{document}
